% Generated by tools/edn_latex.py from reports/edn.md.
% Do not edit: change reports/edn.md and run `make edn`.
\ednentry{
  id = {EDN-32},
  title = {Split proportions 60/10/10/20 and the project seed},
  date = {2026-09-29},
  milestone = {M2: Reproducibility},
  topic = {Evaluation Protocol},
  participants = {@lukas2510},
  decision = {The non-\texttt{ES} listings are split 60\,\% train, 10\,\% validation, 10\,\% calibration and 20\,\% test, grouped by \texttt{seller\_\allowbreak{}group\_\allowbreak{}id}, with the project seed \texttt{20251108} (the snapshot date). Both are pinned in \texttt{params.\allowbreak{}yaml}.},
  rationale = {\emph{Alternatives considered:}
\begin{itemize}[leftmargin=*, topsep=2pt]
\item \textbf{Option A (chosen): 60/10/10/20.}
\newline Pros: keeps the 20\,\% test share that the model card's SC-04 analysis and \texttt{reports/\allowbreak{}analysis/\allowbreak{}make\_\allowbreak{}support\_\allowbreak{}results.\allowbreak{}txt} already assume, so the published evidence stays true without a re-run. Calibration and validation get about 9,800 rows each, which is ample for conformal calibration and for early stopping.
\newline Cons: the smallest training set of the three options, about 58,700 rows.
\item \textbf{Option B: 70/10/10/10.}
\newline Pros: the most training data.
\newline Cons: halves the test set. \texttt{make\_\allowbreak{}support\_\allowbreak{}results.\allowbreak{}txt} gives the test share below which each supported make falls under SC-04's 500-row bar: Volvo at 14.9\,\% and Suzuki at 13.0\,\%. At a 10\,\% test share both drop out, so SC-04 would silently stop checking two of the eleven supported makes and the model card would have to be corrected.
\item \textbf{Option C: 64/8/8/20.}
\newline Pros: the same 20\,\% test share with slightly more training data.
\newline Cons: a less round number to explain in the report, for about 3,900 extra training rows.
\end{itemize}
\par
\emph{Why:} The test share is the binding constraint, not the training share. SC-04 only checks a segment once it holds 500 test rows, so the proportions decide how many makes the quality gate actually covers, and the model card already states which four of the eleven supported makes fall short at 20\,\%. Choosing anything below about 15\,\% would quietly enlarge that set, which is the kind of change that is invisible until someone re-reads the analysis. The seed is the snapshot date rather than a random number so that it is recognisable in a log and obviously not tuned.},
  ai = {Information seeking, Alternative assessment, Recommendation},
  response = {Accepted},
  assessment = {AI found the coupling the ticket did not mention: it read \texttt{make\_\allowbreak{}support\_\allowbreak{}results.\allowbreak{}txt}, noticed that its sensitivity table was computed at an assumed 20\,\% test share while the four-way split was still unpinned, and pointed out that pinning a smaller test share would invalidate a number already published in the model card. It laid out the three options against that constraint and recommended A. Lukas accepted.},
  aievidence = {Claude Code session on 2026-09-29: asked to pin the split, AI reported the per-make sensitivity thresholds from the committed analysis output and framed the three options around the 15\,\% floor they imply.},
  evidence = {\href{https://github.com/mlops-2627q1-mds-upc/MLOPS_Recommenditos/blob/main/params.yaml}{\texttt{params.\allowbreak{}yaml}}; \href{https://github.com/mlops-2627q1-mds-upc/MLOPS_Recommenditos/blob/main/reports/analysis/make_support_results.txt}{reports/analysis/make\_support\_results.txt}; \href{https://github.com/mlops-2627q1-mds-upc/MLOPS_Recommenditos/blob/main/docs/docs/dataset-card.md}{dataset card}; \href{https://github.com/mlops-2627q1-mds-upc/MLOPS_Recommenditos/issues/32}{issue \#32}.},
}
